\uuid{3dHG}
\titre{Composée/DL usuels}
\auteur{hanani ; mzari}
\organisation{exo7}
\licence{CC BY-NC-SA 4.0}
\datecreate{2020-01-20}
\niveau{L1}
\module{Analyse}
\chapitre{Développement limité}
\sousChapitre{Calculs}
\difficulte{3}
\theme{Opérations sur les DL, Composée, DL usuels}
\exo7qcmid{lille-2.developpements-limites.022}

\contenu{
\texte{
Soit $a\in \Rr$ et $f$ une fonction. On dit que $f(x)$ est un $o\left((x-a)^n\right)$ au voisinage de $a$ et on écrit $f(x)=o((x-a)^n)$ si $\displaystyle \lim _{x\to a}\frac{f(x)}{(x-a)^n}=0$. On note par $DL_n(a)f(x)$ le développement limité de $f$ à l'ordre $n$ en $a$. On note aussi $DL_n(a^+)f(x)$ (resp. $DL_n(a^-)f(x)$) le développement limité de $f$ à l'ordre $n$ à droite en $a$ (resp. à gauche en $a$).
}

\begin{question}
\qclassification{125.02}
\qtags{motcle=Composée/DL usuels}
Soit $\displaystyle f(x)=(1+x)^{1/x}$. Parmi les affirmations suivantes, cocher celles qui sont vraies :
\begin{answers}
\good{$\displaystyle DL_2(0)\frac{\ln (1+x)}{x}=1-\frac{x}{2}+\frac{x^2}{3}+o(x^2)$.}
\bad{$\displaystyle DL_2(0)\mathrm{e}^{1+u}=1+u+\frac{u^2}{2}+o(u^2)$.}
\bad{$\displaystyle DL_2(0)f(x)=\frac{5}{2}-2x+\frac{3x^2}{2}+o(x^2)$.}
\good{$\displaystyle DL_2(0)f(x)=\mathrm{e}-\frac{\mathrm{e}}{2}x+\frac{11\mathrm{e}}{24}x^2+o(x^2)$.}
\end{answers}
\begin{explanations}
Le $\displaystyle DL_3(0)\ln (1+x)$ donne : $\displaystyle DL_2(0)\frac{\ln (1+x)}{x}=1-\frac{x}{2}+\frac{x^2}{3}+o(x^2)$.
\vskip0mm
Ensuite, $\mathrm{e}^{1+u}=\mathrm{e}.\mathrm{e}^u=\mathrm{e}.\left(1+u+\frac{u^2}{2}+o(u^2)\right)$. Ainsi, avec $\displaystyle u=-\frac{x}{2}+\frac{x^2}{3}+o(x^2)$, on obtient :
$$f(x)=\mathrm{e}^{1-\frac{x}{2}+\frac{x^2}{3}+o(x^2)}=\mathrm{e}-\frac{\mathrm{e}}{2}x+\frac{11\mathrm{e}}{24}x^2+o(x^2).$$
\end{explanations}
\end{question}
}
